\chapter{Raisonner, définir et démontrer}
\section*{Exercice 1.1 -- Nier sans ambiguïté}
En prenant « positif » au sens $\ge0$, la première négation est : « il existe un indice $n$ tel que $u_n<0$ ». La seconde est : « pour tout entier $n$, $u_n\le3$ ». Nier échange les quantificateurs et remplace $>$ par $\le$, et non par $<$.

\section*{Exercice 1.2 -- Implication et réciproque}
Écrire $n=2k+1$ donne $n^2=4k^2+4k+1=2(2k^2+2k)+1$, impair. Pour la réciproque, sa contraposée suffit : si $n=2k$ est pair, $n^2=2(2k^2)$ est pair. Les deux implications sont donc vraies.

\section*{Exercice 1.3 -- Images réciproques}
Un élément $x$ appartient au membre de gauche si et seulement si $f(x)\in A\cap B$, donc si et seulement si $f(x)\in A$ et $f(x)\in B$. Cela équivaut à $x\in f^{-1}(A)\cap f^{-1}(B)$. Cette chaîne d'équivalences prouve l'égalité des ensembles.

\section*{Exercice 1.4 -- Compter une réunion}
Les $7$ élèves communs sont comptés deux fois dans $18+15$. Le cardinal de la réunion vaut donc $18+15-7=26$. Son complément dans la classe contient $32-26=6$ élèves.

\section*{Exercice 1.5 -- Une récurrence facultative}
Au rang $0$, les deux membres valent $1$. Si $2^n\ge n+1$, alors $2^{n+1}\ge2n+2\ge n+2$ puisque $n\ge0$. L'hérédité et l'initialisation donnent le résultat pour tous les entiers naturels.

\section*{Exercice 1.6 -- Ordre des quantificateurs}
La première est vraie : pour chaque $x$, choisir $y=x+1$. La seconde est fausse : pour n'importe quel $y$ proposé, le choix $x=y$ contredit $y>x$. Dans la première phrase, le choix de $y$ dépend de $x$ ; dans la seconde, un seul $y$ doit fonctionner pour tous.
